\section{Rotor-angle error as an exact coordinate transformation}
\label{sec:angle}
\subsection{Discover the transformation from the axes}
Define the active counterclockwise rotation
\begin{equation}
 \rotationmatrix(\alpha)=\begin{bmatrix}\cos\alpha&-\sin\alpha\\
                  \sin\alpha&\cos\alpha\end{bmatrix},\qquad
 \hat\gamma=\gamma+\angleerror.\label{eq:rotation-definition}
\end{equation}
A positive $\angleerror$ places the estimated d axis counterclockwise from the
true d axis. Coordinates of a fixed physical vector therefore rotate by the
\emph{negative} angle: $\hat\current=\rotationmatrix(-\angleerror)\current$ and
$\hat\voltage=\rotationmatrix(-\angleerror)\voltage$. A constant offset has zero time
derivative, so it leaves electrical speed unchanged.
\begin{figure}[tb]
\centering
\begin{tikzpicture}[scale=1.05]
 \coordinate (O) at (0,0); \coordinate (P) at (2.5,1.8);
 \draw[coordinate vector] (-.4,0)--(3.5,0) node[right]{$d$};
 \draw[coordinate vector] (0,-.4)--(0,3) node[above]{$q$};
 \draw[estimated axis] (0,0)--(20:3.6) node[right]{$\hat d$};
 \draw[estimated axis] (0,0)--(110:2.9) node[above]{$\hat q$};
 \draw[direction arrow,estimated quantity] (.8,0) arc (0:20:.8);
 \node[estimated quantity] at (1.25,.2) {$\angleerror>0$};
 \draw[projection line] (P)--(2.5,0) node[below]{$i_d$};
 \draw[projection line] (P)--(0,1.8) node[left]{$i_q$};
 \coordinate (Dh) at ($(O)!(P)!(20:3.6)$);
 \coordinate (Qh) at ($(O)!(P)!(110:2.9)$);
 \draw[guide pattern,estimated quantity] (P)--(Dh);
 \draw[guide pattern,estimated quantity] (P)--(Qh);
 \node[below right,estimated quantity] at (Dh) {$\hat i_d$};
 \node[above left,estimated quantity] at (Qh) {$\hat i_q$};
 \draw[direction arrow,strong line,neutral quantity] (O)--(P) node[above right]{$\vect i$};
 \node[annotation,text width=55mm] at (5.6,1.5)
 {Same physical vector; different components.\\[3pt]
  $\hat\current=\rotationmatrix(-\angleerror)\current$\\
  $\current=\rotationmatrix(\angleerror)\hat\current$};
\end{tikzpicture}
\caption{True and estimated rotor frames. Positive angle error advances the
estimated axes counterclockwise. Dashed lines show orthogonal component
projections; the angle is exaggerated for visibility.}
\label{fig:frames}
\end{figure}


The constitutive law accepts true current coordinates. Recover them from the
observed ones, evaluate the true map, and project its output back onto the
estimated axes. This gives the exact observed magnetic map
\begin{equation}
 \boxed{\hat\flux_\gamma(\hat\current)=\rotationmatrix(-\angleerror)
       \flux\bigl(\rotationmatrix(\angleerror)\hat\current\bigr).}
 \label{eq:pullback}
\end{equation}
The subscript here excludes reconstruction errors. It is legitimate even
when every stored sample was already transformed using the unknown wrong
angle: the independent variable on the left is the \emph{observed} current.
The true angle is a latent relation between its coordinates and those in the
magnetic model, not a missing measured independent variable.

There are two operations. The inner rotation changes where the magnetics is
evaluated. The outer rotation changes how that flux is resolved into d and q.
Keeping only the first operation models a current-argument error but omits
flux projection. The two can cancel, so omission can even predict a spurious
effect in an isotropic machine.
\begin{figure}[tb]
\centering
\begin{tikzpicture}
 \begin{scope}[xshift=-3.4cm]
  \draw[coordinate vector] (-1.5,0)--(1.8,0) node[right]{$i_d$};
  \draw[coordinate vector] (0,-1.9)--(0,2.2) node[above]{$i_q$};
  \draw[estimated axis] (0,0)--(110:1.9) node[above left]{$+\hat i_q$};
  \draw[estimated axis] (0,0)--(290:1.9) node[below right]{$-\hat i_q$};
  \draw[guide pattern] (110:1.6)--(-.547,0);
  \draw[guide pattern] (290:1.6)--(.547,0);
  \node[operating point,estimated region] at (110:1.6) {};
  \node[operating point,estimated region] at (290:1.6) {};
  \node[compact text,align=center] at (-1.15,-.65) {negative $i_d$\\field weakening};
  \node[compact text,align=center] at (1.3,.75) {positive $i_d$\\magnetization};
  \node[annotation text] at (0,-2.5) {Argument moves: $i_d=-\hat i_q\sin\angleerror$};
 \end{scope}
 \begin{scope}[xshift=3.7cm]
  \draw[coordinate vector] (-.4,0)--(2.4,0) node[right]{$d$};
  \draw[coordinate vector] (0,-.7)--(0,1.8) node[above]{$q$};
  \draw[estimated axis] (0,0)--(20:2.4) node[right]{$\hat d$};
  \draw[estimated axis] (0,0)--(110:1.5) node[above]{$\hat q$};
  \draw[strong line,direction arrow,neutral quantity] (0,0)--(1.9,0) node[below right]{$\flux_{\rm pm}$};
  \coordinate (P) at (1.9,0);
  \coordinate (Q) at ($(0,0)!(P)!(110:1.5)$);
  \draw[guide pattern,estimated quantity] (P)--(Q);
  \draw[strong line,torque quantity,direction arrow] (0,0)--(Q);
  \node[compact text,align=center] at (1,-1.25)
   {PM projection on $\hat q$:\\$-\psi_{\rm pm}\sin\angleerror$};
  \node[annotation text] at (.5,-2.5) {Output is projected onto estimated axes};
 \end{scope}
\end{tikzpicture}

\medskip
\begin{tikzpicture}[node distance=7mm]
 \node[diagnostic block] (i) {Observed current\\$\hat\current$};
 \node[diagnostic block,right=of i] (t) {True current\\$\rotationmatrix(\angleerror)\hat\current$};
 \node[diagnostic block,right=of t] (f) {True magnetic flux\\$\flux(\rotationmatrix(\angleerror)\hat\current)$};
 \node[diagnostic block,right=of f] (h) {Observed flux\\$\rotationmatrix(-\angleerror)\flux(\cdot)$};
 \draw[direction arrow] (i)--(t);\draw[direction arrow] (t)--(f);\draw[direction arrow] (f)--(h);
\end{tikzpicture}
\caption{Two causes of the angle fingerprint. Left: a sweep with
$\hat i_d=0$ follows a tilted line in true current space. Right: the same
PM vector acquires a negative q projection for positive angle error. The
bottom chain describes the full map in observed coordinates, even when the
true angle is unknown. Angles are exaggerated.}
\label{fig:angle-pullback}
\end{figure}


An independent check uses reconstruction: $\rotationgenerator$ commutes with every planar
rotation, so rotating both sides of \cref{eq:vector-voltage} gives the same
voltage balance in the estimated frame. With perfect signals and correct
resistance, its reconstructed flux is exactly \cref{eq:pullback}. With a
resistance mismatch, the complete relation is
\begin{equation}
 \hat\flux(\hat\current)=\hat\flux_\gamma(\hat\current)
                +\frac{\resistanceerror}{\omegae}\rotationgenerator\hat\current.
 \label{eq:combined-exact}
\end{equation}
Under these assumptions the resistance term remains exactly additive even
for finite constant angle offsets. No mixed resistance--angle term is needed
when the map is indexed by observed current.

\subsection{Rotation preserves integrability while rotating the symmetry axis}
The chain rule gives
\begin{equation}
 \hat\flux_\gamma(\hat\current)=
 \nabla_{\hat\current}\coenergy(\rotationmatrix(\angleerror)\hat\current),\qquad
 \hat\Ldiff_\gamma=\rotationmatrix(-\angleerror)\Ldiff(\rotationmatrix(\angleerror)\hat\current)
                         \rotationmatrix(\angleerror).
 \label{eq:rotated-potential}
\end{equation}
This is an orthogonal congruence of a symmetric Hessian and therefore remains
symmetric. A pure constant coherent angle offset has zero curl, even for
finite offset and saturated magnetics. It can violate parity in the chosen
axes because its natural reflection is now
$\rotationmatrix(-\angleerror)\diag(1,-1)\rotationmatrix(\angleerror)$ instead of $\diag(1,-1)$.
Energy consistency alone can fit the rotated magnetic map and does not
identify angle; knowledge of the correct magnetic symmetry axis is essential.

\subsection{The observed zero-d-current sweep}
Before linearizing, inspect the geometry in \cref{fig:angle-pullback}:
\begin{equation}
 \hat i_d=0,\quad\hat i_q=a
 \quad\Longrightarrow\quad
 i_d=-a\sin\angleerror,\quad i_q=a\cos\angleerror.\label{eq:zero-d-sweep}
\end{equation}
For positive offset and positive q current the true d current is negative;
the reflected negative-q point has positive true d current. One half of the
sweep thus enters field weakening and the other additional magnetization.
Even with an even true d-flux map, they are no longer compared at the same
true d current. The subsequent flux projection adds a second contribution.

\subsection{Linearization at a fixed observed current label}
To calculate the local change, expand the inner rotation first:
\begin{align}
 \rotationmatrix(\angleerror)\current&=\current+\angleerror\rotationgenerator\current+\bigO{\angleerror^2},\\
 \flux(\rotationmatrix(\angleerror)\current)&=
 \flux(\current)+\angleerror\Ldiff\rotationgenerator\current+o(|\angleerror|).
\end{align}
Apply $\rotationmatrix(-\angleerror)=\matr I-\angleerror\rotationgenerator+\bigO{\angleerror^2}$ next.
With the numerical observed argument denoted by $\current$ for readability,
\begin{equation}
 \boxed{\delta\flux_\gamma=
 \angleerror\bigl(\Ldiff\rotationgenerator\current-\rotationgenerator\flux\bigr).}
 \label{eq:angle-sensitivity}
\end{equation}
The remainder is $o(|\angleerror|)$ under $C^2$ co-energy and is
$\bigO{\angleerror^2}$ if the flux Jacobian is locally Lipschitz, as it is for
the polynomial example. The two components are
\begin{align}
 \boxed{\delta\psi_{d,\gamma}=
 \angleerror(-L_{dd}i_q+L_{dq}i_d+\psiq),}\label{eq:angle-d}\\
 \boxed{\delta\psi_{q,\gamma}=
 \angleerror(-L_{qd}i_q+L_{qq}i_d-\psid).}\label{eq:angle-q}
\end{align}
The first term in \cref{eq:angle-sensitivity} is motion along the magnetic
surface, the second rotation of its vector output. Both have flux units;
angles are dimensionless. Their difference makes saliency, PM bias, and
local saturation visible. Only first derivatives of flux are required for
this first-order result. Higher derivatives can bound a remainder or generate
higher-order approximations; they are not inputs to this sensitivity.

\subsection{A linear model as a sign and intuition check}
Only now specialize to
\begin{equation}
 \psid=\psi_{\rm pm}+\Ld i_d,\qquad\psiq=\Lq i_q.
\end{equation}
For arbitrary current the sensitivity is
$[(L_q-L_d)i_q,\ (L_q-L_d)i_d-\psi_{\rm pm}]\trans$.
At $i_d=0$ it becomes
\begin{equation}
 \delta\psi_{d,\gamma}=\angleerror(\Lq-\Ld)i_q,\qquad
 \delta\psi_{q,\gamma}=-\angleerror\psi_{\rm pm}.\label{eq:linear-check}
\end{equation}
With $L_d=L_q=L$, isotropic current-generated flux is $L\current$.
Rotating its argument and then rotating its output back cancels exactly.
Only the PM vector rotates: its d change is $\psi_{\rm pm}(\cos\angleerror-1)$,
second order, while its q projection is $-\psi_{\rm pm}\sin\angleerror$.
Thus the q-even channel can be highly sensitive despite having no direct PM
term in the correctly aligned q-flux law.

\paragraph{Saturation changes sensitivity, not the ideal parity.}
At correct alignment saturation alone cannot violate \cref{eq:flux-parity}.
At wrong alignment, operating-point-dependent differential gains change
\cref{eq:angle-sensitivity}. Cross-saturation can reduce a cancellation or
reinforce it, making an offset more visible at some points and less visible
at others. There is no universal monotonic increase with saturation.
For larger offsets evaluate \cref{eq:pullback} directly instead of building
an unnecessary high-order Taylor model.
